\section{弹性体的物理描述}
\label{r2ch07:sec:description}

弹性体（elastic body）在外力撤去后能够回到平衡构形。真实材料还可能具有黏性、塑性和与受力历史有关的迟滞；本章先研究形变小、响应线性、过程无耗散的理想弹性体。对每个材料点，必须同时记录它的参考位置 $\vect R$ 与当前时刻的位置 $\vect r$。运动映射和位移场定义为
\begin{equation}
 \vect r=\vect r(\vect R,t),\qquad
 \vect u(\vect R,t)=\vect r(\vect R,t)-\vect R,\qquad
 \vect u(\vect R+\dd\vect R)-\vect u(\vect R)
 =(\nabla\vect u)\cdot\dd\vect R+O(|\dd\vect R|^2).
 \label{r2ch07:eq:001}
\end{equation}
这里的梯度 $\nabla\vect u$ 是二阶张量。任何二阶张量都可唯一分解为对称部分和反对称部分：
\begin{equation}
 \nabla\vect u=\boldsymbol\varepsilon+\boldsymbol\Omega,\qquad
 \boldsymbol\varepsilon=\frac12\bigl(\nabla\vect u+(\nabla\vect u)^{\mathsf T}\bigr),\qquad
 \boldsymbol\Omega=\frac12\bigl(\nabla\vect u-(\nabla\vect u)^{\mathsf T}\bigr).
 \label{r2ch07:eq:002}
\end{equation}
反对称部分 $\boldsymbol\Omega$ 描述局部小转动；纯平移加纯转动不改变任意两点距离。真正描述小形变的是对称的应变张量（strain tensor）$\boldsymbol\varepsilon$：
\begin{equation}
 \varepsilon_{ij}=\frac12\left(\frac{\partial u_i}{\partial X_j}
 +\frac{\partial u_j}{\partial X_i}\right),\qquad
 \boldsymbol\varepsilon=
 \begin{pmatrix}
 \varepsilon_x&\gamma_{xy}/2&\gamma_{zx}/2\\
 \gamma_{xy}/2&\varepsilon_y&\gamma_{yz}/2\\
 \gamma_{zx}/2&\gamma_{yz}/2&\varepsilon_z
 \end{pmatrix}.
 \label{r2ch07:eq:003}
\end{equation}
对角元是正应变，非对角元的一倍 $\gamma_{ij}=2\varepsilon_{ij}$ 是工程剪应变。图~\ref{r2ch07:fig:001} 把微元先平移、转正，再与虚线原状比较，因而不会把刚体转动误判为剪切：
\begin{equation}
 \varepsilon_x=\frac{\delta L_x}{L_x},\qquad
 \gamma_{xy}=\frac{\partial u_y}{\partial x}+\frac{\partial u_x}{\partial y}
 \simeq\alpha.
 \label{r2ch07:eq:004}
\end{equation}
正应变和剪应变均无量纲，式~\eqref{r2ch07:eq:004} 只在位移梯度远小于 $1$ 时成立。
\figStrain

\subsection{应力张量}
\label{r2ch07:sec:strainstress}

内力不能只用“某一点上的力”描述，因为还必须说明该点处切面的方向。取有向面元 $\dd\vect S=\uvec n\,\dd S$，约定 $\uvec n$ 所指一侧材料对另一侧材料的面力为 $\dd\vect F$。在连续介质极限内，这个关系对面元是线性的：
\begin{equation}
 \dd\vect F=\boldsymbol\sigma\cdot\dd\vect S,\qquad
 \dd F_i=\sum_j\sigma_{ij}\dd S_j.
 \label{r2ch07:eq:005}
\end{equation}
$\boldsymbol\sigma$ 称为应力张量（stress tensor）。取一个边长趋于零的微小立方体，对其力矩平衡，体力矩相对于面积力矩是更高阶小量，因此必须有 $\sigma_{ij}=\sigma_{ji}$。于是
\begin{equation}
 \boldsymbol\sigma=
 \begin{pmatrix}
 \sigma_x&\tau_{xy}&\tau_{zx}\\
 \tau_{xy}&\sigma_y&\tau_{yz}\\
 \tau_{zx}&\tau_{yz}&\sigma_z
 \end{pmatrix},\qquad
 \sigma_n=\frac{\dd F_n}{\dd S},\qquad
 \tau=\frac{\dd F_\tau}{\dd S}.
 \label{r2ch07:eq:006}
\end{equation}
图~\ref{r2ch07:fig:002} 中，正应力垂直切面，剪应力平行切面。这里统一采用“拉应力为正”的固体力学约定，因此静水压强 $p>0$ 对应 $\boldsymbol\sigma=-p\boldsymbol 1$，不能把压强的正号直接当成拉应力正号。
\figStress

\subsection{各向同性线弹性与弹性模量}
\label{r2ch07:sec:constitutive}

一般晶体的应力—应变线性关系需要四阶弹性张量。应力、应变的对称性与无耗散条件会减少独立参数；对各向同性介质，旋转坐标系不能改变材料规律，只剩两个独立常数。最常用的写法是拉梅形式
\begin{equation}
 \boxed{\boldsymbol\sigma=2G\boldsymbol\varepsilon
 +\lambda\operatorname{tr}(\boldsymbol\varepsilon)\boldsymbol 1},
 \label{r2ch07:eq:007}
\end{equation}
其中 $G$ 是剪变模量，$\lambda$ 是另一拉梅参数。分量形式为
\begin{equation}
 \sigma_i=2G\varepsilon_i+\lambda(\varepsilon_x+\varepsilon_y+\varepsilon_z),
 \qquad \tau_{ij}=G\gamma_{ij}.
 \label{r2ch07:eq:008}
\end{equation}
图~\ref{r2ch07:fig:003} 分别表示纯剪切、均匀体积改变和单轴拉伸。体弹性模量 $K$、杨氏模量 $E$、泊松比 $\nu$ 与拉梅参数的关系为
\begin{equation}
 K=\lambda+\frac23G,\qquad
 E=\frac{G(3\lambda+2G)}{\lambda+G},\qquad
 \nu=\frac{\lambda}{2(\lambda+G)},\qquad
 E=2G(1+\nu).
 \label{r2ch07:eq:009}
\end{equation}
单轴拉伸时 $\sigma_x=\sigma$、$\sigma_y=\sigma_z=0$，于是 $\varepsilon_x=\sigma/E$，横向应变为 $\varepsilon_y=\varepsilon_z=-\nu\varepsilon_x$。均匀压缩时 $p=-K\,\Delta V/V$。所有弹性模量的量纲均为压强，国际单位是帕斯卡。稳定的各向同性材料要求 $G>0$、$K>0$，从而 $-1<\nu<1/2$。
\figElasticModuli

\section{弹性棒、弹性绳、弹性膜与三维弹性体}
\label{r2ch07:sec:models}

棒、绳、膜和三维介质的几何维数不同，但建模步骤一致：先由位移梯度求应变，再由本构关系求应力，最后对微元使用牛顿第二定律。相同结构也保证能量能写成局域密度的积分。

\subsection{弹性棒：一维纵波}
\label{r2ch07:subsec:rod}

设均匀直棒横截面积为 $S$、体密度为 $\rho$、杨氏模量为 $E$，$\xi(x,t)$ 是材料点沿棒方向的位移。图~\ref{r2ch07:fig:004} 中原长 $\dd x$ 的微元变为 $\dd x+\dd\xi$，所以
\begin{equation}
 \varepsilon=\frac{\partial\xi}{\partial x},\qquad
 \sigma=E\frac{\partial\xi}{\partial x},\qquad
 \dd F=S\frac{\partial\sigma}{\partial x}\dd x.
 \label{r2ch07:eq:010}
\end{equation}
令 $\dd F=\rho S\dd x\,\partial_t^2\xi$，消去 $S\dd x$ 得
\begin{equation}
 \boxed{\frac{\partial^2\xi}{\partial t^2}
 =\frac E\rho\frac{\partial^2\xi}{\partial x^2}},
 \qquad v_{\rm rod}=\sqrt{\frac E\rho}.
 \label{r2ch07:eq:011}
\end{equation}
该模型忽略横向惯性、截面翘曲和波长可与棒直径相比时的色散，只适用于细长棒的小振幅纵振动。单位体积的动能、势能和总能量密度为
\begin{equation}
 \mathcal T=\frac12\rho\left(\frac{\partial\xi}{\partial t}\right)^2,\qquad
 \mathcal V=\frac12E\left(\frac{\partial\xi}{\partial x}\right)^2,\qquad
 \mathcal E=\mathcal T+\mathcal V.
 \label{r2ch07:eq:012}
\end{equation}
\figRodModel

\subsection{张紧弦：一维传播、二维偏振}
\label{r2ch07:subsec:string}

质量线密度为 $\mu$ 的细绳预先由恒定张力 $T$ 拉紧。设两个横向位移为 $\xi(x,t),\eta(x,t)$。小斜率条件 $|\partial_x\xi|,|\partial_x\eta|\ll1$ 下，图~\ref{r2ch07:fig:005} 中两端张力的横向分量差给出
\begin{equation}
 \frac{\partial^2\xi}{\partial t^2}=\frac T\mu\frac{\partial^2\xi}{\partial x^2},\qquad
 \frac{\partial^2\eta}{\partial t^2}=\frac T\mu\frac{\partial^2\eta}{\partial x^2},\qquad
 v_{\rm string}=\sqrt{\frac T\mu}.
 \label{r2ch07:eq:013}
\end{equation}
两个方程相互独立，所以弦波可线偏振、椭圆偏振或圆偏振。单位长度能量密度为
\begin{equation}
 \mathcal T=\frac\mu2(\xi_t^2+\eta_t^2),\qquad
 \mathcal V=\frac T2(\xi_x^2+\eta_x^2),\qquad
 V-TL_0\simeq\int\mathcal V\,\dd x.
 \label{r2ch07:eq:014}
\end{equation}
最后一式来自 $V=T\int\dd s$ 的小斜率展开，说明恒张力弦的横振势能等于额外弧长乘张力。
\figStringModel

\subsection{张紧膜：二维传播}
\label{r2ch07:subsec:membrane}

设膜的质量面密度为 $\mu_s$、各向同性张力系数为 $\Sigma$，横位移为 $\zeta(x,y,t)$。图~\ref{r2ch07:fig:006} 中任意闭合边界上的面内张力，利用格林公式化为区域内的恢复力，得到
\begin{equation}
 \boxed{\frac{\partial^2\zeta}{\partial t^2}
 =\frac\Sigma{\mu_s}\left(\frac{\partial^2\zeta}{\partial x^2}
 +\frac{\partial^2\zeta}{\partial y^2}\right)},\qquad
 v_{\rm membrane}=\sqrt{\frac\Sigma{\mu_s}}.
 \label{r2ch07:eq:015}
\end{equation}
单位面积能量密度为
\begin{equation}
 \mathcal T=\frac{\mu_s}{2}\zeta_t^2,\qquad
 \mathcal V=\frac\Sigma2(\zeta_x^2+\zeta_y^2),\qquad
 V-\Sigma S_0\simeq\iint\mathcal V\,\dd x\dd y.
 \label{r2ch07:eq:016}
\end{equation}
这里仍要求斜率小；若膜本身的弯曲刚度不可忽略，方程会出现更高阶空间导数。
\figMembraneModel

\subsection{三维弹性体}
\label{r2ch07:subsec:body}

对体积区域 $D$，边界面力的合力是 $\oint_{\partial D}\boldsymbol\sigma\cdot\dd\vect S$。用高斯定理把它化成体积分，单位体积的内力为 $\nabla\cdot\boldsymbol\sigma$。与式~\eqref{r2ch07:eq:007} 联立，得到纳维—柯西方程
\begin{equation}
 \boxed{\rho\frac{\partial^2\vect u}{\partial t^2}
 =(\lambda+2G)\nabla(\nabla\cdot\vect u)
 -G\nabla\times(\nabla\times\vect u)}.
 \label{r2ch07:eq:017}
\end{equation}
式中使用了 $\nabla^2\vect u=\nabla(\nabla\cdot\vect u)-\nabla\times(\nabla\times\vect u)$。若还存在体力密度 $\vect f$，只需在右端加 $\vect f$。动能与弹性势能密度可写成坐标无关形式
\begin{equation}
 \mathcal T=\frac12\rho|\dot{\vect u}|^2,\qquad
 \mathcal V=\frac12\boldsymbol\sigma:\boldsymbol\varepsilon
 =\frac\lambda2(\nabla\cdot\vect u)^2+G\boldsymbol\varepsilon:\boldsymbol\varepsilon.
 \label{r2ch07:eq:018}
\end{equation}
双点乘定义为 $\mathsf A:\mathsf B=\sum_{ij}A_{ij}B_{ij}$。$G>0,K>0$ 正是保证任意非零应变都有正势能的条件。

\section{弹性波}
\label{r2ch07:sec:waves}

\subsection{一维波动方程的两种求解方法}
\label{r2ch07:subsec:oned}

一维波动方程
\begin{equation}
 \xi_{tt}=v^2\xi_{xx}
 \label{r2ch07:eq:019}
\end{equation}
是线性齐次偏微分方程，解满足叠加原理。分离变量 $\xi=X(x)Q(t)$ 后，两边只能等于同一常数。选择振荡时间因子 $Q=e^{i\omega t}$，空间因子为 $X=e^{\mp ikx}$，且 $\omega=vk$。另一种做法令 $x_\pm=x\pm vt$，方程变成 $\partial^2\xi/(\partial x_+\partial x_-)=0$，直接积分得到达朗贝尔通解
\begin{equation}
 \boxed{\xi(x,t)=f(x-vt)+g(x+vt)}.
 \label{r2ch07:eq:020}
\end{equation}
$f$ 保持波形向 $+x$ 传播，$g$ 向 $-x$ 传播。理想棒、绳模型中 $v$ 与频率无关，故无色散；真实介质在波长接近微观尺度或横向尺度时会偏离这一近似。

\subsection{多维波、传播方向与偏振}
\label{r2ch07:subsec:multidim}

二维膜的简谐平面波及三维矢量平面波可统一写成
\begin{equation}
 \zeta=Ae^{i(\omega t-\vect k\cdot\vect r)},\qquad
 \vect u=\vect A e^{i(\omega t-\vect k\cdot\vect r)},\qquad
 \omega=v|\vect k|.
 \label{r2ch07:eq:021}
\end{equation}
$|\vect k|$ 决定波长，$\vect k/|\vect k|$ 决定传播方向，复振幅的模和幅角决定强度与相位；矢量振幅的方向还给出偏振。

\subsection{三维弹性体中的P波与S波}
\label{r2ch07:subsec:bodywaves}

斯托克斯—亥姆霍兹分解把足够光滑的位移场分成无散部分和无旋部分：
\begin{equation}
 \vect u=\vect u_S+\vect u_P,\qquad
 \nabla\cdot\vect u_S=0,\qquad
 \nabla\times\vect u_P=0.
 \label{r2ch07:eq:022}
\end{equation}
把两部分分别代入式~\eqref{r2ch07:eq:017}。无散部分满足
\begin{equation}
 \vect u_{S,tt}=v_S^2\nabla^2\vect u_S,\qquad
 v_S=\sqrt{\frac G\rho},\qquad
 \vect k\cdot\vect A_S=0.
 \label{r2ch07:eq:023}
\end{equation}
它是位移垂直传播方向的剪切波（S波），有两个独立横向偏振。无旋部分满足
\begin{equation}
 \vect u_{P,tt}=v_P^2\nabla^2\vect u_P,\qquad
 v_P=\sqrt{\frac{\lambda+2G}{\rho}}
 =\sqrt{\frac{K+4G/3}{\rho}},\qquad
 \vect k\times\vect A_P=0.
 \label{r2ch07:eq:024}
\end{equation}
它是位移平行传播方向的压缩波（P波）。图~\ref{r2ch07:fig:007} 明确给出了位移箭头与 $\vect k$ 的相对方向。因为 $K,G>0$，总有 $v_P>v_S$。液体的静态剪变模量近似为零，不能传播体S波；这一事实可用于判断地球外核的液态性质。
\figBodyWaves

液体或气体中 $G\simeq0$，声速为 $v=\sqrt{K/\rho}$。理想气体中的声振动很快，气团来不及与外界换热，应使用绝热关系 $pV^\gamma=\text{常量}$，于是
\begin{equation}
 K=-V\frac{\dd p}{\dd V}=\gamma p,\qquad
 \boxed{v_{\rm sound}=\sqrt{\frac{\gamma p}{\rho}}
 =\sqrt{\frac{\gamma RT}{M_{\rm mol}}}}.
 \label{r2ch07:eq:025}
\end{equation}
量纲检查：$K/\rho$ 的单位为 $(\mathrm{N,m^{-2}})/(\mathrm{kg,m^{-3}})=\mathrm{m^2s^{-2}}$，开方后确为速度。
